% ===================================================================
%  ਸਾਰਣੀ ਨੰ. 11 — ਸ਼ਹਿਰ ਤੇ ਉਸ ਦੀਆਂ ਇਮਾਰਤਾਂ। — ਅੱਗ।
%  Traduction 2026 d'apres texto/io, controlee sur texto/fr, texto/en
%  et texto/hi. Consigne : voir texto/pa/10-sarni-01.tex.
% ===================================================================

%%K t11-apar-1 apar
\VUsaut{2.0mm}
\VUcentre{13.2pt}{90}{ਤੀਜੀ ਲੜੀ}
\VUsaut{3.0mm}
\VUfilet{16mm}
\VUsaut{5.6mm}
\VUtitre{15.0pt}{60}{ਸਾਰਣੀ ਨੰ. 11}
\VUsaut{1.4mm}
\VUpk{11.4pt}{40}{ਸ਼ਹਿਰ ਤੇ ਉਸ ਦੀਆਂ ਇਮਾਰਤਾਂ। --- ਅੱਗ।}
\VUsaut{2.2mm}

%%K t11-tit-1 sub
\VUpk{10.2pt}{40}{ਬਾਹਰਲੀ ਬਸਤੀ।}
\VUsaut{3.0mm}

%%K t11-c1-01-1 p
1. --- ਮੈਂ ਮਹਾਂਸਾਗਰ ਤੋਂ ਸੌ ਕਿਲੋਮੀਟਰ ਦੂਰ ਇਕ ਵੱਡੇ ਸ਼ਹਿਰ ਵਿਚ ਰਹਿੰਦਾ ਹਾਂ, ਜੋ ਫਿਰ ਵੀ
ਪਹਿਲੇ ਦਰਜੇ ਦੀ \VUgras{ਬੰਦਰਗਾਹ} \textsuperscript{(1)} ਹੈ। ਜੋ ਦਰਿਆ ਉਸ ਨਾਲ ਵਗਦਾ ਹੈ
ਉਹ ਚਾਰ ਸੌ ਮੀਟਰ ਚੌੜਾ ਹੈ ਤੇ ਇਕ ਬਹੁਤ ਸੋਹਣਾ ਲੰਗਰਗਾਹ ਬਣਾਉਂਦਾ ਹੈ।

%%K t11-c1-02-1 p
2. --- \VUgras{ਘਾਟਾਂ} \textsuperscript{(2)} ਨਾਲ ਦਾ ਸਾਰਾ ਹਿੱਸਾ ਬਿਲਕੁਲ ਸਮੁੰਦਰੀ ਤੇ
ਸਨਅਤੀ ਰੂਪ ਰੱਖਦਾ ਹੈ, ਤੇ ਇਕ ਇਹੋ ਜਿਹੀ \VUgras{ਬਾਹਰਲੀ ਬਸਤੀ} \textsuperscript{(3)}
ਬਣਾਉਂਦਾ ਹੈ ਜਿਸ ਵਿਚ ਸਿਰਫ਼ ਮਲਾਹ ਤੇ ਮਜ਼ਦੂਰ ਰਹਿੰਦੇ ਹਨ। ਇਹ \VUgras{ਕਾਰਖ਼ਾਨਿਆਂ}
\textsuperscript{(4)} ਵਿਚ ਤੇ \VUgras{ਗੈਸ-ਘਰ} \textsuperscript{(5)} ਵਿਚ ਕੰਮ ਕਰਦੇ ਹਨ,
ਜਿਸ ਦੀ ਉੱਚੀ \VUgras{ਚਿਮਨੀ} \textsuperscript{(6)} ਤੇ \VUgras{ਗੈਸ-ਟੈਂਕੀ} ਦੂਰੋਂ ਹੀ
ਵਿਖਾਈ ਦਿੰਦੀ ਹੈ।

%%K t11-c1-03-1 p
3. --- ਵਿਚਕਾਰਲੇ ਜੁਗ ਵਿਚ ਸ਼ਹਿਰ ਕਿਲ੍ਹਾਬੰਦ ਸੀ, ਤੇ ਉਸ ਦੀਆਂ ਫ਼ਸੀਲਾਂ ਦੇ ਕੁਝ ਨਿਸ਼ਾਨ ਹੁਣ
ਵੀ ਲੱਭਦੇ ਹਨ, ਖ਼ਾਸ ਕਰ ਕੇ ਉਹ \VUgras{ਬੂਹਾ} \textsuperscript{(8)} ਜੋ ਪੁਰਾਣੇ
\VUgras{ਗੜ੍ਹ-ਕਿਲ੍ਹੇ} \textsuperscript{(7)} ਦਾ ਹੈ, ਤੇ ਜਿਸ ਦੇ ਦੋਹਾਂ ਪਾਸੇ ਉਸ ਦੇ ਦੋ
\VUgras{ਮੁਨਾਰੇ} \textsuperscript{(9)} ਹਨ, ਜੋ ਹੁਣ \VUgras{ਕੈਦਖ਼ਾਨੇ} ਦੇ ਕੰਮ ਆਉਂਦੇ ਹਨ।

%%K t11-c1-04-1 p
4. --- ਇਸ ਸਾਰੇ \VUgras{ਮੁਹੱਲੇ} ਵਿਚ, ਜੋ ਬਹੁਤ ਪੁਰਾਣਾ ਲਗਦਾ ਹੈ, ਸਿਰਫ਼ ਇਕ ਇਮਾਰਤ ਰਤਾ
ਨਵੀਂ ਹੈ : \VUgras{ਚੁੰਗੀ-ਘਰ} \textsuperscript{(10)}। ਉਹ ਘਾਟ ਤੇ ਉਸ ਨਿੱਕੀ
\VUgras{ਗਲੀ} \textsuperscript{(11)} ਵਿਚਕਾਰ ਖੜਾ ਹੈ ਜੋ ਪੁਰਾਣੇ \VUgras{ਸ਼ਹਿਰੀ ਘਰ}
\textsuperscript{(12)} ਤਕ ਜਾਂਦੀ ਹੈ। ਉਸ ਆਖ਼ਰੀ ਇਮਾਰਤ ਨੂੰ ਉਸ ਵੱਡੇ
\VUgras{ਘੰਟਾ-ਮੁਨਾਰੇ} \textsuperscript{(13)} ਤੋਂ ਸੌਖਾ ਹੀ ਪਛਾਣਿਆ ਜਾ ਸਕਦਾ ਹੈ ਜੋ ਪੁਰਾਣੇ
ਸ਼ਹਿਰ ਉੱਤੇ ਛਾਇਆ ਰਹਿੰਦਾ ਹੈ।

%%K t11-c1-05-1 p
5. --- ਗਲੀ ਦੇ ਦੂਜੇ ਪਾਸੇ ਚੁੰਗੀ-ਘਰ ਦੀ ਹੀ ਵਜ਼ਾ ਦੀ ਇਕ ਇਮਾਰਤ \VUgras{ਖੜੀ} ਹੈ : ਉਹ
\VUgras{ਵਪਾਰ-ਘਰ} \textsuperscript{(14)} ਹੈ। ਉਸ ਦਾ ਇਕ ਮੂੰਹ ਇਕ ਨਿੱਕੇ ਚੌਕ ਵੱਲ
ਖੁੱਲ੍ਹਦਾ ਹੈ ਜਿੱਥੇ \VUgras{ਪ੍ਰਬੰਧਕੀ ਘਰ} \textsuperscript{(15)} ਹੈ। ਸਾਡਾ ਸ਼ਹਿਰ, ਅਸਲ
ਵਿਚ, ਰਾਜਧਾਨੀ ਨਾ ਹੁੰਦੇ ਹੋਏ ਵੀ ਇਕ ਅਹਿਮ ਜ਼ਿਲ੍ਹਾ-ਸਦਰ ਮੁਕਾਮ ਹੈ।

%%K t11-tit-2 sub
\VUsaut{5.0mm}
\VUpk{10.2pt}{40}{ਸ਼ਹਿਰ ਦਾ ਉੱਪਰਲਾ ਹਿੱਸਾ।}
\VUsaut{3.4mm}

%%K t11-c1-06-1 p
6. --- ਛਾਉਣੀ, ਜੋ ਕਾਫ਼ੀ ਵੱਡੀ ਹੈ, ਵੱਡੀਆਂ ਤੇ ਬਹੁਤ ਸਿਹਤਮੰਦ \VUgras{ਬੈਰਕਾਂ}
\textsuperscript{(16)} ਵਿਚ ਰਹਿੰਦੀ ਹੈ; ਉਹ \VUgras{ਸ਼ਹਿਰੀ ਬਾਗ਼} \textsuperscript{(17)}
ਦੇ ਜੰਗਲੇ ਤਕ ਫੈਲੀਆਂ ਹਨ। ਜੋ \VUgras{ਚੌੜੀ ਸੜਕ} \textsuperscript{(19)} ਇਸ ਬਾਗ਼ ਤਕ ਜਾਂਦੀ
ਹੈ, ਉਹ \VUgras{ਬੱਚਤ ਬੈਂਕ} \textsuperscript{(18)} ਦੇ ਪਿੱਛਿਓਂ ਲੰਘ ਕੇ
\VUgras{ਵੱਡੇ ਗਿਰਜੇ} \textsuperscript{(20)} ਦੇ ਚੌਕ ਉੱਤੇ ਨਿਕਲਦੀ ਹੈ।

%%K t11-c1-07-1 p
7. --- ਇਹ ਸਾਰੀਆਂ ਇਮਾਰਤਾਂ ਇਕ ਨਿੱਕੀ ਟਿੱਬੀ ਉੱਤੇ ਪੌੜੀਆਂ ਵਾਂਗ ਚੜ੍ਹਦੀਆਂ ਹਨ, ਜਿੱਥੇ ਸਾਡਾ
ਵੱਡਾ \VUgras{ਹਾਈ ਸਕੂਲ} \textsuperscript{(21)} ਵੀ ਹੈ।

%%K t11-c1-07-2 p
ਜੇ ਤੁਸੀਂ ਉਸ ਹੌਲੀ ਢਾਲ ਵਾਲੀ ਸੜਕ ਤੋਂ ਉਤਰੋ ਜੋ ਵੱਡੀ ਸੜਕ ਨਾਲ ਮਿਲਦੀ ਹੈ, ਤਾਂ ਤੁਸੀਂ
\VUgras{ਅਦਾਲਤ} \textsuperscript{(22)} ਤਕ ਪਹੁੰਚਦੇ ਹੋ, ਜਿਸ ਦੇ ਸਾਹਮਣੇ ਇਕ ਸੋਹਣਾ
\VUgras{ਸਰਬਜਨਕ ਬਾਗ਼} \textsuperscript{(26)} ਹੈ। ਇਹ \VUgras{ਚੌਕ} ਬਹੁਤ ਵੱਡਾ ਨਹੀਂ, ਪਰ
ਸ਼ਹਿਰ ਦੇ ਵਿਚਕਾਰ ਉਹ ਹਰਿਆਵਲ ਤੇ ਠੰਢ ਦਾ ਨਖ਼ਲਿਸਤਾਨ ਬਣਾਉਂਦਾ ਹੈ। ਸੋ ਸ਼ਹਿਰ ਵਾਲੇ ਉਸ ਵਿਚ
ਖ਼ੂਬ ਆਉਂਦੇ ਹਨ, ਜਿਨ੍ਹਾਂ ਨੂੰ \VUgras{ਕਹਿਵਾਘਰ} \textsuperscript{(25)} ਦੇ ਛੱਜੇ ਦੇ ਹੇਠਾਂ
ਬਹਿ ਕੇ ਉਸ ਫ਼ੌਜੀ ਵਾਜੇ ਨੂੰ ਸੁਣਨਾ ਚੰਗਾ ਲਗਦਾ ਹੈ ਜੋ ਵੀਰਵਾਰ ਤੇ ਐਤਵਾਰ ਨੂੰ ਲਾਨ ਤੇ
ਕਿਆਰੀਆਂ ਵਿਚਕਾਰ ਬਣੇ \VUgras{ਵਾਜੇ ਦੇ ਮੰਚ} \textsuperscript{(27)} ਵਿਚ ਵੱਜਦਾ ਹੈ।

%%K t11-c1-08-1 p
8. --- ਉਨ੍ਹਾਂ ਹੀ ਲਾਨਾਂ ਵਿਚੋਂ ਇਕ ਉੱਤੇ ਕਾਂਸੀ ਦੀ ਇਕ \VUgras{ਮੂਰਤੀ}
\textsuperscript{(28)} ਖੜੀ ਹੈ, ਸਾਡੇ ਇਕ ਸ਼ਹਿਰ ਵਾਲੇ ਦੀ ਯਾਦ ਵਿਚ ਚੁੱਕੀ ਯਾਦਗਾਰ, ਜਿਸ ਨੇ
ਆਪਣੀ ਜਨਮ-ਭੋਇੰ ਦੀ ਵੱਡੀ ਸੇਵਾ ਕੀਤੀ।

%%K t11-tit-3 sub
\VUsaut{4.0mm}
\VUpk{10.2pt}{40}{ਆਉਣਾ-ਜਾਣਾ।}
\VUsaut{3.4mm}

%%K t11-c1-09-1 p
9. --- ਸਾਡਾ ਸ਼ਹਿਰ ਹਾਲੇ ਬਿਜਲੀ ਨਾਲ ਰੌਸ਼ਨ ਨਹੀਂ; ਉਸ ਵਿਚ ਸਿਰਫ਼ \VUgras{ਗੈਸ ਦੇ ਲੈਂਪ}
\textsuperscript{(29)} ਹਨ। ਪਰ \VUgras{ਸਥਾਨਕ ਅਖ਼ਬਾਰ} ਇਸ ਗੱਲ ਲਈ ਮੁਹਿੰਮ ਚਲਾ ਰਹੇ ਹਨ ਕਿ
ਰੌਸ਼ਨੀ ਦਾ ਉਹ ਪ੍ਰਬੰਧ ਸੁਧਾਰਿਆ ਜਾਵੇ, ਜਿਵੇਂ ਟਰਾਮਾਂ ਦਾ \VUgras{ਖਿੱਚਣ ਦਾ ਪ੍ਰਬੰਧ}
ਪਹਿਲਾਂ ਹੀ ਸੁਧਾਰਿਆ ਜਾ ਚੁੱਕਾ ਹੈ। \VUgras{ਘੋੜਿਆਂ} ਨਾਲ \VUgras{ਖਿੱਚੀਆਂ} ਜਾਣ ਵਾਲੀਆਂ
ਪੁਰਾਣੀਆਂ ਗੱਡੀਆਂ ਦੀ ਥਾਂ ਹੁਣ ਕੰਮ ਦੀਆਂ \VUgras{ਬਿਜਲੀ ਦੀਆਂ ਟਰਾਮਾਂ}
\textsuperscript{(31)} ਆ ਗਈਆਂ ਹਨ, ਜੋ ਬਹੁਤ \VUgras{ਘੱਟ} ਕਿਰਾਏ ਵਿਚ ਮੁਸਾਫ਼ਰਾਂ ਨੂੰ
\VUgras{ਸ਼ਹਿਰ ਦੇ} ਇਕ ਸਿਰੇ ਤੋਂ \VUgras{ਦੂਜੇ ਤਕ} ਛੇਤੀ ਪਹੁੰਚਾ ਦਿੰਦੀਆਂ ਹਨ।

%%K t11-c1-10-1 p
10. --- ਅਦਾਲਤ ਕੋਲ \VUgras{ਬੈਂਕ} \textsuperscript{(32)} ਹੈ, ਜਿੱਥੇ ਵਪਾਰੀ ਹੁੰਡੀਆਂ
ਭੰਨਾਈਆਂ ਜਾਂਦੀਆਂ ਹਨ। \VUgras{ਬੈਂਕ ਦੇ ਹਰਕਾਰੇ} \textsuperscript{(33)} ਲੋਕਾਂ ਦੇ
ਘਰ-ਘਰ ਜਾ ਕੇ ਬਕਾਇਆ ਵਸੂਲਦੇ ਹਨ।

%%K t11-tit-4 sub
\VUsaut{4.0mm}
\VUpk{10.2pt}{40}{ਕਲਾ ਤੇ ਵਿੱਦਿਆ।}
\VUsaut{3.4mm}

%%K t11-c1-11-1 p
11. --- ਕੋਲ ਦੀ ਉਹ ਸੋਹਣੀ ਇਮਾਰਤ \VUgras{ਅਜਾਇਬ-ਘਰ} ਹੈ। ਉਸ ਵਿਚ ਇਕੋ ਵੇਲੇ ਸ਼ਹਿਰ ਦਾ
\VUgras{ਕਿਤਾਬ-ਘਰ} \textsuperscript{(34)}, ਕੁਝ ਸੋਹਣੇ
\VUgras{ਕੁਦਰਤੀ ਇਤਿਹਾਸ ਦੇ ਸੰਗ੍ਰਹਿ} \textsuperscript{(35)} (ਇਕ ਕੁਦਰਤੀ ਇਤਿਹਾਸ ਦਾ
ਅਜਾਇਬ-ਘਰ) ਤੇ ਇਕ ਸੋਹਣਾ \VUgras{ਤਸਵੀਰ-ਘਰ} \textsuperscript{(36)} ਹੈ, ਜੋ ਇਕ ਸ਼ਾਨਦਾਰ
ਪੱਕੀ \VUgras{ਨੁਮਾਇਸ਼} ਬਣਾਉਂਦਾ ਹੈ।

%%K t11-c1-11-2 p
ਉਹ ਹਫ਼ਤੇ ਵਿਚ ਦੋ ਵਾਰ ਲੋਕਾਂ ਲਈ ਖੁੱਲ੍ਹਦਾ ਹੈ, ਪਰ ਵੇਖਣ ਵਾਲੇ ਕਿਸੇ ਵੀ ਦਿਨ ਉਸ ਨੂੰ ਵੇਖ
ਸਕਦੇ ਹਨ, \VUgras{ਰਾਖੇ} \textsuperscript{(37)} ਨੂੰ ਥੋੜੀ ਬਖ਼ਸ਼ੀਸ਼ ਦੇ ਕੇ।
\VUgras{ਮੁਸੱਵਰ} \textsuperscript{(38)} ਉੱਥੇ ਕੰਮ ਕਰਨ ਆਉਂਦੇ ਹਨ। ਉਨ੍ਹਾਂ ਨੂੰ ਆਪਣੀਆਂ
ਤਿਪਾਈਆਂ ਦੇ \VUgras{ਸਾਹਮਣੇ} ਵੇਖਿਆ ਜਾ ਸਕਦਾ ਹੈ, ਹੱਥ ਵਿਚ \VUgras{ਰੰਗ-ਪੱਟੀ} ਚੁੱਕੀ,
ਮਸ਼ਹੂਰ ਤਸਵੀਰਾਂ ਦੀ ਨਕਲ ਉਤਾਰਦੇ।

%%K t11-c1-12-1 p
12. --- ਅਜਾਇਬ-ਘਰ ਦੇ ਪਿਛਲੀਆਂ ਉਚਾਈਆਂ ਉੱਤੇ ਉਹ \VUgras{ਗੁੰਬਦ} \textsuperscript{(40)}
ਜੋ \VUgras{ਹਸਪਤਾਲ} \textsuperscript{(39)} ਦਾ ਹੈ, ਤੇ ਉਸ
\VUgras{ਅਧਿਆਪਕਾਂ ਦੇ ਸਿਖਲਾਈ ਸਕੂਲ} \textsuperscript{(41)} ਦੀਆਂ ਇਮਾਰਤਾਂ ਮੁਸ਼ਕਲ ਨਾਲ
ਵਿਖਾਈ ਦਿੰਦੀਆਂ ਹਨ ਜਿੱਥੇ ਸਾਡੇ ਸ਼ਹਿਰੀ ਸਕੂਲਾਂ ਦੇ ਅਧਿਆਪਕ ਸਿਖਲਾਈ ਪਾਏ ਸਨ। ਰੰਗ-ਮੰਚ ਦੇ
ਚੌਕ ਉੱਤੇ ਇਕ \VUgras{ਡਾਕਘਰ} \textsuperscript{(43)} ਹੈ, ਜੋ ਇਕ \VUgras{ਨਿੱਜੀ ਮਕਾਨ}
\textsuperscript{(42)} ਵਿਚ ਬੈਠਾ ਹੈ।

%%K t11-tit-5 sub
\VUsaut{5.0mm}
\VUpk{11.4pt}{40}{ਅੱਗ।}
\VUsaut{3.0mm}

%%K t11-tit-6 sub
\VUpk{10.2pt}{40}{ਅੱਗ ਦੀ ਹਾਕ।}
\VUsaut{3.4mm}

%%K t11-c2-01-1 p
1. --- ਸ਼ਹਿਰ ਵਿਚ ਡਰਾਉਣੀ ਖ਼ਬਰ ਫੈਲ ਰਹੀ ਹੈ : ਰੰਗ-ਮੰਚ ਨੂੰ ਹੁਣੇ ਹੁਣੇ ਅੱਗ ਲੱਗ ਗਈ! ਜਦੋਂ
ਤਕ ਮੈਂ \VUgras{ਚੌਕ} \textsuperscript{(96)} ਪਹੁੰਚਦਾ ਹਾਂ, ਭੀੜ \VUgras{ਪਟੜੀ}
\textsuperscript{(46)} ਤੇ \VUgras{ਸੜਕ} \textsuperscript{(47)} ਉੱਤੇ ਜੁੜ ਚੁੱਕੀ ਹੈ।
\VUgras{ਡਾਕ ਦੇ ਬਾਬੂ} \textsuperscript{(44)} ਤੇ \VUgras{ਡਾਕੀਏ}
\textsuperscript{(45)} ਡਾਕਘਰ ਤੋਂ ਬਾਹਰ ਆ ਰਹੇ ਹਨ। ਇਕ \VUgras{ਟਰਾਮ ਚਾਲਕ}
\textsuperscript{(48)} ਨੂੰ ਆਪਣੀ ਟਰਾਮ ਰੋਕਣੀ ਪਈ ਹੈ ਤਾਂ ਜੋ ਸ਼ਹਿਰ ਦੀ ਇਕ
\VUgras{ਮਰੀਜ਼-ਗੱਡੀ} \textsuperscript{(50)} ਨਿਕਲ ਸਕੇ, ਜੋ ਇਕ \VUgras{ਜਰਾਹ}
\textsuperscript{(52)} ਤੇ ਇਕ \VUgras{ਨਰਸ} \textsuperscript{(51)} ਨੂੰ ਲੈ ਕੇ ਆ ਰਹੀ ਹੈ,
ਇਕ \VUgras{ਜ਼ਖ਼ਮੀ} \textsuperscript{(53)} ਦੀ ਸੰਭਾਲ ਲਈ, ਜਿਸ ਨੂੰ \VUgras{ਸਟਰੈਚਰ}
\textsuperscript{(54)} ਉੱਤੇ \VUgras{ਅਖ਼ਬਾਰ ਦੇ ਖੋਖੇ} \textsuperscript{(55)} ਕੋਲ ਲਿਆਂਦਾ
ਗਿਆ ਹੈ।

%%K t11-c2-02-1 p
2. --- ਕਵਾਇਦ ਤੋਂ ਮੁੜਦੀ ਪੈਦਲ ਫ਼ੌਜ ਦੀ ਇਕ ਟੋਲੀ ਰੁਕ ਗਈ ਹੈ, ਤਾਂ ਜੋ ਘੋੜਿਆਂ ਉੱਤੇ ਸਵਾਰ
\VUgras{ਸ਼ਹਿਰੀ ਰਾਖਿਆਂ} \textsuperscript{(61)} ਨਾਲ ਪ੍ਰਬੰਧ ਬਣਾਈ ਰੱਖਣ ਵਿਚ ਮਦਦ ਕਰੇ।
\VUgras{ਘੋੜਸਵਾਰਾਂ} ਦੇ ਘੋੜੇ ਦੋ ਪੈਰਾਂ ਉੱਤੇ ਖੜੇ ਹੋ ਜਾਂਦੇ ਹਨ; ਉਹ \VUgras{ਸਿਪਾਹੀਆਂ}
\textsuperscript{(57)} ਜਾਂ \VUgras{ਜ਼ਾਂਦਾਰਮਾਂ} \textsuperscript{(63)} ਨਾਲੋਂ ਚੰਗੀ
ਤਰ੍ਹਾਂ \VUgras{ਤਮਾਸ਼ਬੀਨਾਂ} \textsuperscript{(56)} ਨੂੰ ਪਿੱਛੇ ਧੱਕ ਸਕਦੇ ਹਨ। ਜ਼ਾਂਦਾਰਮ
ਵੈਸੇ ਵੀ ਬਹੁਤ ਰੁੱਝੇ ਹਨ। ਉਨ੍ਹਾਂ ਵਿਚੋਂ ਇਕ \VUgras{ਪੁਲਿਸ ਦੇ ਸਿਪਾਹੀਆਂ}
\textsuperscript{(64)} ਨੂੰ ਦੱਸ ਰਿਹਾ ਹੈ ਕਿ ਇਕ ਹੋਰ \VUgras{ਜ਼ਖ਼ਮੀ}
\textsuperscript{(65)} ਨੂੰ ਕਿੱਥੇ ਲੈ ਜਾਣਾ ਹੈ, ਇਕ ਸੂਰਮਾ ਅੱਗ ਬੁਝਾਊ ਵਾਲਾ ਜੋ ਪੌੜੀ ਤੋਂ
ਡਿੱਗ ਪਿਆ ਹੈ।

%%K t11-c2-03-1 p
3. --- \VUgras{ਅਫ਼ਸਰ} \textsuperscript{(66)} ਠੀਕ ਠੀਕ ਦੱਸਦਾ ਹੈ ਕਿ ਜੋ
\VUgras{ਭਾਫ਼ ਦਾ ਪੰਪ} \textsuperscript{(67)} ਹੁਣੇ ਆ ਰਿਹਾ ਹੈ ਉਸ ਨੂੰ ਕਿੱਥੇ ਰੱਖਣਾ ਹੈ।
ਉਸ ਤਕੜੇ ਜੰਤਰ ਦੀ ਉਡੀਕ ਵਿਚ ਉਸ ਨੇ ਇਕ \VUgras{ਹੱਥ ਦਾ ਪੰਪ} \textsuperscript{(68)}
ਲਗਵਾ ਦਿੱਤਾ ਹੈ, ਜਿਸ ਦੀ ਲੰਮੀ \VUgras{ਨਾਲੀ} \textsuperscript{(69)} ਅੱਗ ਉੱਤੇ ਪਾਣੀ ਦੀਆਂ
ਧਾਰਾਂ ਡੋਲ੍ਹ ਰਹੀ ਹੈ। ਛੇ ਅੱਗ ਬੁਝਾਊ ਵਾਲੇ ਉਸ ਨੂੰ ਚਲਾ ਰਹੇ ਹਨ, ਜਦੋਂ ਤਕ ਉਨ੍ਹਾਂ ਦਾ ਸਾਥੀ
\VUgras{ਟੂਟੀ} \textsuperscript{(70)} ਫੜੀ \VUgras{ਧਾਰ} \textsuperscript{(71)} ਨੂੰ ਅੱਗ ਦੇ
ਵਿਚਕਾਰ ਸਾਧਦਾ ਹੈ।

%%K t11-c2-04-1 p
4. --- ਰੰਗ-ਮੰਚ ਦੇ ਦੂਜੇ ਪਾਸੇ ਵੱਡੀ \VUgras{ਪੌੜੀ} \textsuperscript{(72)} ਖੜੀ ਕੀਤੀ ਗਈ
ਹੈ। ਉਹ ਅੱਗ ਬੁਝਾਊ ਵਾਲਿਆਂ ਨੂੰ ਉਨ੍ਹਾਂ ਲਾਟਾਂ ਕੋਲ ਪਹੁੰਚਣ ਦਿੰਦੀ ਹੈ ਜੋ ਕਾਲੇ
\VUgras{ਧੂੰਏਂ} \textsuperscript{(75)} ਦੇ ਵਰੋਲਿਆਂ ਵਿਚਕਾਰ ਉੱਠਦੀਆਂ ਹਨ।

%%K t11-tit-7 sub
\VUsaut{5.0mm}
\VUpk{10.2pt}{40}{ਸੂਰਮਗਤੀ ਦੇ ਕੰਮ।}
\VUsaut{3.4mm}

%%K t11-c2-05-1 p
5. --- \VUgras{ਅੱਗ} \textsuperscript{(74)} \VUgras{ਰੰਗ-ਮੰਚ} \textsuperscript{(73)}
ਦੇ ਡਿਉੜ੍ਹੀ-ਕਮਰੇ ਵਿਚ ਲੱਗੀ, ਉਸ \VUgras{ਓਪੇਰਾ} ਦੀ ਰਿਹਰਸਲ ਦੇ ਵੇਲੇ ਜਿਸ ਦਾ ਪਹਿਲਾ
\VUgras{ਖੇਲ} ਕੱਲ੍ਹ ਹੋਣ ਵਾਲਾ ਸੀ। ਖ਼ੁਸ਼ਕਿਸਮਤੀ ਨਾਲ ਹਾਲ ਵਿਚ ਸਿਰਫ਼ ਕੁਝ
\VUgras{ਵੇਖਣ ਵਾਲੇ} \textsuperscript{(76)} ਸਨ। ਵਿਚਾਰਿਆਂ ਨੇ \VUgras{ਬਾਲਾਖ਼ਾਨੇ}
\textsuperscript{(77)} ਵਿਚ ਪਨਾਹ ਲਈ, ਜਿੱਥੇ ਸੂਰਮੇ \VUgras{ਅੱਗ ਬੁਝਾਊ ਵਾਲੇ}
\textsuperscript{(86)} ਪੌੜੀ ਤੇ ਰੱਸੀਆਂ ਨਾਲ ਉਨ੍ਹਾਂ ਦੀ ਮਦਦ ਨੂੰ ਆ ਰਹੇ ਹਨ।

%%K t11-c2-06-1 p
6. --- \VUgras{ਡਿਉੜ੍ਹੀ} \textsuperscript{(78)} ਦੇ ਹੇਠਾਂ, ਜਿੱਥੇ
\VUgras{ਨਾਟਕ ਦੀ ਸੂਚਨਾ} \textsuperscript{(79)} ਖੇਲ ਦਾ ਐਲਾਨ ਕਰਦੀ ਹੈ,
\VUgras{ਸਾਜ਼ ਵਾਲੇ} \textsuperscript{(81)} ਨੱਸਦੇ ਵਿਖਾਈ ਦਿੰਦੇ ਹਨ, ਜੋ
\VUgras{ਸਾਜ਼-ਮੰਡਲੀ} \textsuperscript{(80)} ਦੇ ਹਨ, ਆਪਣੇ ਸਾਜ਼ ਚੁੱਕੀ :
\VUgras{ਵਾਇਲਨ} \textsuperscript{(81)}, \VUgras{ਬੰਸਰੀ} \textsuperscript{(82)},
\VUgras{ਚੈਲੋ} \textsuperscript{(83)}, \VUgras{ਟਰੋਂਬੋਨ} \textsuperscript{(84)},
\VUgras{ਢੋਲ} \textsuperscript{(85)}, ਵਗ਼ੈਰਾ। \VUgras{ਸਾਜ਼-ਸਾਮਾਨ}
\textsuperscript{(87)} ਦਾ ਇਕ ਹਿੱਸਾ ਬਚਾਉਣ ਦਾ ਯਤਨ ਕੀਤਾ ਜਾ ਰਿਹਾ ਹੈ।

%%K t11-c2-07-1 p
7. --- \VUgras{ਅਦਾਕਾਰਾਂ} \textsuperscript{(89)} ਤੇ \VUgras{ਅਦਾਕਾਰਾਵਾਂ}
\textsuperscript{(88)} ਨੂੰ, ਗਵੱਈਆਂ ਤੇ ਗਵੱਈਆਂ ਨੂੰ, ਆਪਣੇ ਮੰਚ ਦੀਆਂ \VUgras{ਪੁਸ਼ਾਕਾਂ}
\textsuperscript{(90)} ਵਿਚ ਹੀ ਨੱਸਣਾ ਪਿਆ ਹੈ। ਮੁੱਖ ਗਵੱਈਆ ਇਕ \VUgras{ਪੱਤਰਕਾਰ}
\textsuperscript{(92)} ਨੂੰ ਸਮਝਾ ਰਿਹਾ ਹੈ ਕਿ ਇਹ ਕਿਵੇਂ ਹੋਇਆ। \VUgras{ਖ਼ਬਰਨਵੀਸ} ਪਹਿਲੇ
ਹੀ ਪਲ ਤੋਂ ਅਫ਼ਸਰਾਂ ਨਾਲ ਦੌੜਿਆ ਆਇਆ ਸੀ : \VUgras{ਪ੍ਰਬੰਧਕ} \textsuperscript{(91)},
\VUgras{ਸ਼ਹਿਰ-ਮੁਖੀ} \textsuperscript{(93)}, \VUgras{ਪੁਲਿਸ ਅਫ਼ਸਰ}
\textsuperscript{(94)} ਤੇ ਇਕ \VUgras{ਜੱਜ} \textsuperscript{(95)}, ਜੋ ਪੜਤਾਲ ਕਰਨਗੇ ਕਿ
ਜ਼ਿੰਮੇਵਾਰੀ ਕਿਸ ਦੀ ਹੈ।
\VUsaut{6.0mm}
\VUfilet{22mm}
